\documentclass[11pt]{article}
\usepackage[T1]{fontenc} % classnotes.sty omits this; without it, a literal
                          % ">" in regular text renders as an upside-down
                          % question mark. cis3300-accessible.sty already
                          % loads T1, worth adding to classnotes.sty too.
\usepackage[margin=1in]{geometry}
\usepackage[pagenumbers]{classnotes}

% This document is Java throughout except two data-file listings, which
% override this locally with language={}. showstringspaces=false is added
% locally because the shared classnotes.sty "sh" style does not set it
% (cis3300-accessible.sty's style does); worth adding upstream too.
\lstset{language=Java, basicstyle=\footnotesize\ttfamily, showstringspaces=false,
  aboveskip=5pt, belowskip=5pt,
  morecomment=[l][\color{codecmt}\itshape]{//}}

\classnotesmeta{CIS 3300 - Objects, Files, and ArrayLists: The Racer Class}

\begin{document}

\classnotestitle{CIS 3300 \textbullet\ In-Class Demo}{Objects, Files, and ArrayLists}{Reading three file formats, and writing one back out}

\noindent\textbf{\color{amber}A word before you read.}\enspace
  These are my notes from over the years, cleaned up so they are readable. A few things to
  keep in mind:
 
  \begin{itemize}[leftmargin=1.4em]
    \item \textbf{They are not a study guide, and they are not gospel.} Not everything here
      appears in a lecture, and not every topic has notes yet; I will write more as we go.
    \item \textbf{I write from memory and revise every semester.} A typo or an imprecise
      statement can survive longer than it should.
    \item \textbf{I use AI to help} convert documents to markdown, check for issues, and add
      explanations. I do not trust it completely, so I am still editing these myself.
    \item \textbf{Java changes.} The basics we cover are pretty resilient, but some features
      do move over time.
  \end{itemize}
 
\noindent So: if you find something that looks wrong, contradicts the book, or just
does not behave the way I claim when you run it, tell me. You will be doing the other students
a favor, and you will understand the material better for having caught it. When in doubt, the
textbook is the tiebreaker, and the compiler beats us both.

A quick note before starting: today's demo uses a small custom class,
\texttt{Racer}, to hold a few pieces of data together. This is meant
as a light, practical introduction, just enough to build and use
objects in this week's lab. Classes and objects get covered in full
depth in a later class, so do not worry about understanding everything
about how or why a class works yet, only how to create one and call
its methods.

With that said, today's goal is to see the same \texttt{Racer}
objects handled several different ways: built by hand, read from a
plain space-separated text file, read from a CSV file two different
ways, and finally written back out to a file. The \texttt{Racer}
class and its \texttt{ArrayList<Racer>} stay the same throughout;
only where the data comes from, or goes to, changes.

\section{Project Setup}

Do this once, before Part 1, and everything below will just work.

\begin{enumerate}[itemsep=2pt]
  \item Create a new Eclipse project.
  \item Right click the \texttt{src} folder, choose
    \textbf{New > Package}, and name it \texttt{racer}. Every class
    today goes in this package, starting with \texttt{package racer;}
    as its first line.
  \item Right click the project name, choose \textbf{New > Folder},
    and name it \texttt{files}. Every data file you read today, and
    everything your code writes out, lives here.
\end{enumerate}

\section{Part 1: Building Racer Objects by Hand}

A class is a blueprint. It describes what a \texttt{Racer} object
knows about itself (its name, its vehicle, and its top speed) but it
does not create any actual racers on its own. Each time the program
writes \texttt{new Racer(...)}, one new object gets built from that
blueprint and stored in memory.

\subsection{The Racer class}

\lstinputlisting{Racer.java}

The constructor is the piece that runs at the exact moment
\texttt{new Racer(...)} executes. Its only job is to take the values
passed in and store them in the object's fields, using
\texttt{this.name = name;} to tell Java which \texttt{name} is the
field and which is the parameter.

\subsection{Creating racers and putting them in a list}

\lstinputlisting{RacerDemoPart1.java}

Every call to \texttt{racers.add(new Racer(...))} builds one object
and adds it to the list in a single step. Once the objects are in the
\texttt{ArrayList}, \texttt{racers.get(i)} hands back one
\texttt{Racer} object at a time, and calling \texttt{.getName()} or
\texttt{.getTopSpeed()} on it reads that object's fields. The
fastest-racer loop at the bottom does not use anything special to
search the list, it just walks every element with a regular for loop
and keeps track of the best one seen so far, an idea that comes back
later today.

\section{Part 2: Reading a Space-Separated File}

Not every data file is a CSV. \texttt{lapTimes.txt} below has one
racer per line, with a name and a lap time separated by a space
instead of a comma.

\subsection{lapTimes.txt}

\lstinputlisting[language={}]{lapTimes.txt}

\subsection{Reading lapTimes.txt}

\lstinputlisting{RacerDemoPart2.java}

\texttt{Scanner} already knows how to split on whitespace, so no
manual splitting is needed here. Calling \texttt{.next()} reads the
next word, and \texttt{.nextDouble()} reads the next number, each one
stopping automatically at the next space or line break. This is the
same idea you will use to read a file with several numbers per line,
separated by spaces, later this week.

\section{Part 3: Reading a CSV File Without a Library}

\texttt{racers.csv} below is the same three racers as Part 1, this
time saved to a file with commas between the fields.

\subsection{racers.csv}

\lstinputlisting[language={}]{racers.csv}

\subsection{Reading racers.csv by hand}

\lstinputlisting{RacerDemoPart3.java}

This reads each line as one long \texttt{String} and then calls
\texttt{line.split(",")} to cut it into pieces wherever a comma
appears. It works here because nothing in \texttt{racers.csv}
contains a comma inside a field. If a vehicle name were written as
\texttt{"Kart, Deluxe"}, with a comma of its own, \texttt{split(",")}
would cut that name into two pieces and shift every value after it
over by one column, silently producing the wrong data instead of an
error. Real-world CSV data is often messier than today's example,
which is why a dedicated library is worth using.

\section{Part 4: Reading the Same CSV File with OpenCSV}

This part reads the exact same \texttt{racers.csv} from Part 3 again,
using the OpenCSV library instead of \texttt{split(",")}, so you can
compare the two approaches directly.

\subsection{Adding OpenCSV and Apache Commons Lang to your Eclipse project}

OpenCSV needs one small helper library, Apache Commons Lang, to run.
The steps below are a simple way to add both jars to a project by
hand. This is not necessarily the most robust way to manage a library
in a larger project, but it is enough for today, and it will be
walked through live in class.

\begin{enumerate}[itemsep=2pt]
  \item Download OpenCSV from
    \url{https://sourceforge.net/projects/opencsv/}, and Apache
    Commons Lang from
    \url{https://commons.apache.org/proper/commons-lang/download_lang.cgi}.
    Each download is an archive; extract both and find the plain
    \texttt{.jar} file in each one, not the \texttt{-sources} or
    \texttt{-javadoc} jar.
  \item In Eclipse, right click your project in the Package Explorer,
    choose \textbf{New > Folder}, and name it \texttt{lib}. Copy both
    jar files into that folder.
  \item Right click the project name, choose
    \textbf{Build Path > Configure Build Path}, open the
    \textbf{Libraries} tab, choose \textbf{Classpath}, then
    \textbf{Add JARs} and select both jars in your \texttt{lib}
    folder.
  \item Click \textbf{Apply and Close}. The line
    \texttt{import com.opencsv.CSVReader;} should no longer show an
    error.
\end{enumerate}

% Instructor note, not shown to students: this simplified two-jar
% setup (OpenCSV + Apache Commons Lang) is enough for the basic
% CSVReader usage shown in this demo. OpenCSV's full formal dependency
% chain also includes commons-text, commons-beanutils,
% commons-collections, and commons-collections4; those are not needed
% for anything shown today, but would be if a feature reaching further
% into OpenCSV were used later. Recheck OpenCSV's current dependency
% list before reusing this approach in a future semester.

\subsection{Reading racers.csv with OpenCSV}

\lstinputlisting{RacerDemoPart4.java}

\texttt{reader.readNext()} returns one line at a time as a String
array, already correctly split at the commas, handling a comma inside
a quoted field the way \texttt{split(",")} could not. It returns
\texttt{null} once the file runs out, which is what the \texttt{while}
loop's condition is checking.

\section{Part 5: Writing Data to a File}

Reading data is only half the picture. This part builds the same
racer list from Part 1 again and writes a short report out to a new
file, \texttt{raceReport.txt}.

\subsection{Writing raceReport.txt}

\lstinputlisting{RacerDemoPart5.java}

\texttt{PrintWriter} has the same \texttt{println(...)} method you
have already used with \texttt{System.out}, except its output goes to
a file instead of the console. The call to \texttt{writer.close()} at
the end is not optional: output can stay held in memory rather than
actually written to disk until the writer is closed, so skipping it
can leave the file empty or cut off partway through.

\section{What to Notice}

\begin{itemize}[itemsep=2pt]
  \item The \texttt{Racer} class did not change across any of the
    five parts above. Objects do not care whether the values used to
    build them were typed directly into the code, read from a file,
    or about to be written back out to one.
  \item \texttt{Scanner}'s \texttt{.next()}, \texttt{.nextInt()}, and
    \texttt{.nextDouble()} methods split on whitespace automatically,
    which is exactly what a space-separated file needs.
  \item \texttt{split(",")} is fine for simple, well-behaved CSV data,
    but it breaks the moment a field contains a comma of its own.
    OpenCSV exists to handle that case correctly.
  \item Writing to a file with \texttt{PrintWriter} looks almost
    identical to printing to the console, the output just goes
    somewhere else, and the file has to be closed when you are done.
\end{itemize}

\notebox{Looking ahead.}{The lab that pairs with this demo asks you to
combine two files into one list, matching values by an ID and updating
objects that already exist in the list rather than only adding new
ones. Everything above shows how to read and write each kind of file
on its own; combining and updating come next.}

\end{document}
